% Required packages: graphicx, booktabs.
% Uses the manuscript's existing natbib configuration.
\section{Experiments}
\label{sec:experiments}

We evaluate Conv on Qwen3-8B and Llama-3.1-8B using LongBench and RULER.
LongBench results are reported for all 16 evaluated tasks, with their
unweighted mean. RULER results are averaged over 13 subtasks at 32K, 64K,
and 128K tokens. We compare with full attention,
MInference~\citep{jiang2024minference},
FlexPrefill~\citep{lai2025flexprefill}, and the XAttention implementation
denoted as Xattn~\citep{xu2025xattention}. Conv and Xattn use a score-sampling
stride of 8. Evaluation details are provided in
Appendix~\ref{app:experimental_details}.

\subsection{Overall Results and Efficiency}
\label{sec:main_results_efficiency}

\subsubsection{Final Benchmark Results}
\label{sec:final_benchmark_results}

Tables~\ref{tab:longbench_llama_topk}--\ref{tab:longbench_qwen_topk}
report the complete LongBench results. Both Conv and Xattn use a retained-block
ratio of 0.70 for the Llama top-$k$ comparison and 0.75 for the Qwen3
comparison. Under these matched budgets, Conv improves the average from
34.10 to 35.22 on Llama and from 44.78 to 44.91 on Qwen3. The gains vary
across tasks: on Llama, improvements on Qasper and TREC coexist with lower
scores on MultiFieldQA-en and MuSiQue. On Qwen3, Conv improves 7 of 16 tasks,
including 2WikiMQA and HotpotQA, but the aggregate improvement is small.

The Llama top-$p=0.90$ comparison gives a different operating point:
Conv scores 36.67 on average, compared with 37.65 for Xattn. Conv performs
better on 4 of 16 tasks, including LCC and RepoBench-P, but has a lower
overall average. Thus, its top-$k$ accuracy gains should not be interpreted
as universal accuracy improvements under top-$p$ selection. Equal
cumulative-score thresholds do not impose equal retained-block budgets.

\begin{table}[t]
    \centering
    \caption{LongBench results on Llama-3.1-8B. Conv and Xattn both use top-$k=0.70$. Full attention, MInference, and FlexPrefill are reference configurations; the top-$k$ setting applies only to Conv and Xattn. Higher is better.}
    \label{tab:longbench_llama_topk}
    \small
    \setlength{\tabcolsep}{5pt}
    \begin{tabular}{lrrrrr}
        \toprule
        Task & Full & MInference & FlexPrefill & Xattn & Conv \\
        \midrule
        2WikiMQA & 17.27 & 17.28 & 16.93 & 16.79 & 15.06 \\
        GovReport & 34.71 & 34.45 & 33.59 & 33.32 & 34.32 \\
        HotpotQA & 20.21 & 20.43 & 19.59 & 21.06 & 22.15 \\
        LCC & 54.12 & 54.01 & 53.28 & 44.66 & 44.95 \\
        LSHT & 46.00 & 46.00 & 44.00 & 28.00 & 34.50 \\
        MultiNews & 27.87 & 28.01 & 28.20 & 27.89 & 28.16 \\
        MultiFieldQA-en & 29.33 & 30.74 & 28.17 & 36.42 & 31.82 \\
        MuSiQue & 13.24 & 12.66 & 11.66 & 12.59 & 10.25 \\
        NarrativeQA & 31.53 & 31.77 & 29.73 & 27.82 & 28.61 \\
        Qasper & 25.84 & 25.03 & 26.26 & 25.83 & 31.42 \\
        QMSum & 23.33 & 23.46 & 23.54 & 22.33 & 22.82 \\
        RepoBench-P & 51.50 & 52.04 & 55.07 & 48.02 & 48.58 \\
        SAMSum & 43.80 & 43.71 & 43.36 & 41.65 & 43.12 \\
        TREC & 72.50 & 72.50 & 70.50 & 61.50 & 66.50 \\
        TriviaQA & 91.65 & 91.18 & 88.83 & 82.29 & 85.71 \\
        VCSum & 15.97 & 16.04 & 16.68 & 15.46 & 15.56 \\
        \midrule
        Average & 37.43 & 37.46 & 36.84 & 34.10 & 35.22 \\
        \bottomrule
    \end{tabular}
\end{table}

\begin{table}[t]
    \centering
    \caption{LongBench results on Llama-3.1-8B with top-$p=0.90$ for Conv and Xattn. The three reference-method columns repeat the same results as Table~\ref{tab:longbench_llama_topk}; they are not separate top-$p$ runs. Higher is better.}
    \label{tab:longbench_llama_topp}
    \small
    \setlength{\tabcolsep}{5pt}
    \begin{tabular}{lrrrrr}
        \toprule
        Task & Full & MInference & FlexPrefill & Xattn & Conv \\
        \midrule
        2WikiMQA & 17.27 & 17.28 & 16.93 & 16.22 & 14.65 \\
        GovReport & 34.71 & 34.45 & 33.59 & 34.96 & 34.42 \\
        HotpotQA & 20.21 & 20.43 & 19.59 & 21.20 & 20.22 \\
        LCC & 54.12 & 54.01 & 53.28 & 54.90 & 55.83 \\
        LSHT & 46.00 & 46.00 & 44.00 & 43.50 & 42.00 \\
        MultiNews & 27.87 & 28.01 & 28.20 & 28.09 & 27.90 \\
        MultiFieldQA-en & 29.33 & 30.74 & 28.17 & 30.01 & 28.27 \\
        MuSiQue & 13.24 & 12.66 & 11.66 & 13.25 & 11.46 \\
        NarrativeQA & 31.53 & 31.77 & 29.73 & 30.47 & 27.10 \\
        Qasper & 25.84 & 25.03 & 26.26 & 27.63 & 26.63 \\
        QMSum & 23.33 & 23.46 & 23.54 & 23.32 & 23.01 \\
        RepoBench-P & 51.50 & 52.04 & 55.07 & 54.41 & 56.19 \\
        SAMSum & 43.80 & 43.71 & 43.36 & 44.28 & 42.86 \\
        TREC & 72.50 & 72.50 & 70.50 & 71.50 & 72.00 \\
        TriviaQA & 91.65 & 91.18 & 88.83 & 92.36 & 87.70 \\
        VCSum & 15.97 & 16.04 & 16.68 & 16.25 & 16.49 \\
        \midrule
        Average & 37.43 & 37.46 & 36.84 & 37.65 & 36.67 \\
        \bottomrule
    \end{tabular}
\end{table}

\begin{table}[t]
    \centering
    \caption{LongBench results on Qwen3-8B. Conv and Xattn both use top-$k=0.75$ at evaluation time. This is an inference setting, irrespective of the training-ratio notation in the checkpoint name. Higher is better.}
    \label{tab:longbench_qwen_topk}
    \small
    \setlength{\tabcolsep}{5pt}
    \begin{tabular}{lrrrrr}
        \toprule
        Task & Full & MInference & FlexPrefill & Xattn & Conv \\
        \midrule
        2WikiMQA & 42.95 & 43.04 & 37.26 & 41.00 & 43.80 \\
        GovReport & 33.57 & 33.56 & 33.37 & 33.57 & 33.46 \\
        HotpotQA & 59.34 & 59.06 & 57.85 & 56.71 & 59.18 \\
        LCC & 57.77 & 57.89 & 58.16 & 52.05 & 51.79 \\
        LSHT & 47.50 & 47.50 & 39.75 & 49.50 & 48.50 \\
        MultiNews & 24.76 & 24.92 & 24.66 & 24.62 & 24.93 \\
        MultiFieldQA-en & 54.17 & 53.23 & 53.92 & 54.57 & 53.33 \\
        MuSiQue & 36.51 & 35.81 & 32.62 & 34.63 & 34.04 \\
        NarrativeQA & 26.31 & 26.08 & 27.57 & 27.33 & 26.39 \\
        Qasper & 48.34 & 48.33 & 43.72 & 44.95 & 45.41 \\
        QMSum & 24.09 & 24.02 & 24.44 & 23.86 & 23.99 \\
        RepoBench-P & 56.46 & 56.96 & 58.36 & 55.16 & 54.94 \\
        SAMSum & 44.17 & 44.44 & 44.98 & 43.41 & 43.34 \\
        TREC & 71.50 & 72.00 & 74.00 & 72.50 & 71.50 \\
        TriviaQA & 90.71 & 90.71 & 88.97 & 88.42 & 89.70 \\
        VCSum & 14.06 & 14.17 & 14.27 & 14.19 & 14.22 \\
        \midrule
        Average & 45.76 & 45.73 & 44.62 & 44.78 & 44.91 \\
        \bottomrule
    \end{tabular}
\end{table}

Table~\ref{tab:ruler_main_results} separately reports RULER results. Conv
and Xattn use matched top-$k$ ratios of 0.65 on Qwen3 and 0.70 on Llama.
On Qwen3, the RULER-trained Conv scores 75.89 on average, compared with
75.99 for Xattn: small gains at 32K and 64K are offset by a lower score
at 128K. On Llama, Conv improves the score by
0.97, 2.39, and 1.99 points at 32K, 64K, and 128K, respectively, increasing
the three-length average from 77.83 to 79.61. MInference and FlexPrefill
achieve higher aggregate scores in their respective configurations; these
comparisons are not matched in retained-block budget.

% Qwen Conv uses stage3_extend_96k128k_t065_s8_yarn4_bf16_ema_step9750.
\begin{table}[t]
    \centering
    \caption{RULER results with matched top-$k$ ratios for Conv and Xattn:
        0.65 on Qwen3-8B and 0.70 on Llama-3.1-8B.
        Each length averages 13 subtasks; Avg. is the mean across lengths.
        Other methods use their benchmark configurations rather than
        this fixed keep ratio. ``--'' indicates an unavailable result.
        Higher is better.}
    \label{tab:ruler_main_results}
    \small
    \begin{tabular}{llrrrr}
        \toprule
        Model & Method & 32K & 64K & 128K & Avg. \\
        \midrule
        Qwen3-8B & Full attention & 92.27 & 72.41 & 70.38 & 78.36 \\
        & MInference & 92.58 & 73.60 & 70.17 & 78.78 \\
        & FlexPrefill & 92.33 & 70.90 & 67.92 & 77.05 \\
        & Xattn & 90.65 & 68.92 & 68.39 & 75.99 \\
        & Conv & 90.68 & 69.32 & 67.67 & 75.89 \\
        \addlinespace
        Llama-3.1-8B & Full attention & -- & -- & -- & -- \\
        & MInference & 94.10 & 89.49 & 80.48 & 88.02 \\
        & FlexPrefill & 92.57 & 86.39 & 79.63 & 86.20 \\
        & Xattn & 74.64 & 81.28 & 77.56 & 77.83 \\
        & Conv & 75.62 & 83.67 & 79.55 & 79.61 \\
        \bottomrule
    \end{tabular}
\end{table}

\subsubsection{Fixed-Budget Top-$k$ Comparison}
\label{sec:topk_comparison}

Figure~\ref{fig:topk_tradeoff} presents the Llama top-$k$ sweep, with
numerical results in Table~\ref{tab:topk_tradeoff_values}. Across 32K, 64K, and 128K contexts and keep ratios
from 0.60 to 0.75, Conv improves RULER accuracy in all 12 matched-budget
settings. Gains range from 0.97 to 7.64 points and average 3.07 points.
The reported speedup factors differ by at most $0.02\times$, indicating
similar measured efficiency at the same retained-block budget.

The sweep also identifies cases where Conv matches or exceeds Xattn with
fewer blocks. At 32K, Conv at 0.70 scores 75.62, compared with 75.23 for
Xattn at 0.75. At 128K, Conv at 0.70 scores 79.55, exceeding Xattn at 0.75
by 1.02 points. These cases demonstrate a better accuracy--budget trade-off,
although this ordering does not hold for every pair of adjacent keep ratios.

\begin{figure}[t]
    \centering
    \includegraphics[width=\linewidth]{Figs/ruler3map.png}
    \caption{Llama-3.1-8B top-$k$ comparison at 32K, 64K, and 128K.
        Solid lines and circles show RULER scores; dashed lines and squares
        show speedup over full attention. Numerical results are reported
        in Table~\ref{tab:topk_tradeoff_values}.}
    \label{fig:topk_tradeoff}
\end{figure}

\begin{table}[t]
    \centering
    \caption{Numerical results corresponding to Figure~\ref{fig:topk_tradeoff}.
        Ratio denotes the retained-block fraction. Scores and speedups are
        rounded to two decimal places.}
    \label{tab:topk_tradeoff_values}
    \small
    \setlength{\tabcolsep}{6pt}
    \renewcommand{\arraystretch}{1.1}
    \begin{tabular}{@{}rrrcrr@{}}
        \toprule
        & \multicolumn{2}{c}{RULER score} && \multicolumn{2}{c}{Speedup ($\times$)} \\
        \cmidrule(lr){2-3}\cmidrule(lr){5-6}
        Ratio & Xattn & Conv && Xattn & Conv \\
        \midrule
        \multicolumn{6}{c}{32K context} \\
        0.60 & 65.34 & 69.06 &  & 1.28 & 1.30 \\
        0.65 & 69.63 & 72.51 &  & 1.19 & 1.20 \\
        0.70 & 74.64 & 75.62 &  & 1.12 & 1.12 \\
        0.75 & 75.23 & 82.87 &  & 1.06 & 1.06 \\
        \addlinespace
        \multicolumn{6}{c}{64K context} \\
        0.60 & 72.38 & 76.44 &  & 1.34 & 1.36 \\
        0.65 & 77.02 & 80.78 &  & 1.25 & 1.27 \\
        0.70 & 81.28 & 83.67 &  & 1.18 & 1.19 \\
        0.75 & 83.90 & 86.10 &  & 1.10 & 1.12 \\
        \addlinespace
        \multicolumn{6}{c}{128K context} \\
        0.60 & 72.11 & 74.92 &  & 1.39 & 1.41 \\
        0.65 & 74.20 & 75.97 &  & 1.29 & 1.31 \\
        0.70 & 77.56 & 79.55 &  & 1.21 & 1.22 \\
        0.75 & 78.53 & 81.18 &  & 1.14 & 1.15 \\
        \bottomrule
    \end{tabular}
\end{table}

Fixed top-$k$ selection isolates block-selection quality by controlling
the retained budget. In practical use, however, the budget needed to
preserve useful information can vary across inputs and tasks. We therefore
also evaluate top-$p$ selection, which adapts the retained count according
to cumulative estimated importance. This complements the fixed-budget
comparison, while requiring accuracy and speed to be assessed together.

\subsubsection{Adaptive Top-$p$ Attention Efficiency}
\label{sec:topp_efficiency}

Figure~\ref{fig:topp_speedup} reports Llama attention-path speedups with
top-$p=0.90$ for Conv and Xattn. The benchmark times the methods on shared
Q/K/V tensors; it does not measure end-to-end language-model generation.
Conv achieves $1.76\times$, $2.61\times$, and $3.67\times$ speedup over
full attention at 32K, 64K, and 128K, compared with $1.42\times$,
$2.06\times$, and $2.87\times$ for Xattn. At 128K, the ratio of the reported
speedups is $1.28\times$, corresponding to approximately 22\% less measured
attention time. At 4K and 8K, Conv is slower than both full attention and
Xattn, showing that the benefit depends on context length.

These efficiency results should be read alongside the top-$p$ LongBench
comparison: Conv's mean score is 0.98 points below Xattn at the same
threshold. The measurements therefore establish a trade-off between
accuracy and attention runtime, rather than acceleration at matched
accuracy. Because the timing and accuracy evaluations use different
workloads, we do not interpret the timing ratios as LongBench end-to-end
speedups.

\begin{figure}[t]
    \centering
    \includegraphics[width=0.9\linewidth]{Figs/topp_efficiency_llama.png}
    \caption{Attention-path speedup on Llama-3.1-8B. Conv and Xattn use
        top-$p=0.90$; the other methods use their efficiency-benchmark
        settings. Full attention is normalized to $1.00\times$ separately
        at each length. Values below one indicate slower execution.
        This measures the attention path rather than full-model latency.}
    \label{fig:topp_speedup}
\end{figure}

\subsection{Ablation Studies}
\label{sec:ablations}

\subsubsection{Effect of RULER-Oriented Training}
\label{sec:training_ablation}

We first test whether the learned convolutional refinement benefits from the
RULER-oriented training stage. Table~\ref{tab:ruler_training_ablation}
compares the original $7\times7$ convolution with the RULER-trained variant on
Qwen3-8B, holding the top-$k$ ratio at 0.65. Training improves the RULER score
by 1.30 points at 32K and 0.75 points at 64K. The 128K score changes by
$-0.07$ points, while the average across the three lengths increases from
75.23 to 75.89. Thus, the training stage produces a measurable overall gain,
although the improvement is concentrated at 32K and 64K rather than being
uniform across all lengths.

\begin{table}[t]
    \centering
    \caption{
        Ablation of RULER-oriented training on Qwen3-8B with a $7\times7$
        kernel and top-$k=0.65$. Higher is better.
    }
    \label{tab:ruler_training_ablation}
    \small
    \begin{tabular}{lrrrr}
        \toprule
        Variant & 32K & 64K & 128K & Avg. \\
        \midrule
        Original convolution & 89.38 & 68.57 & \textbf{67.74} & 75.23 \\
        RULER-trained convolution & \textbf{90.68} & \textbf{69.32} & 67.67 & \textbf{75.89} \\
        \midrule
        Difference & +1.30 & +0.75 & $-0.07$ & +0.66 \\
        \bottomrule
    \end{tabular}
\end{table}

% --------------------------------------------------------------------------
% Kernel-size ablation: keep this block commented until the experiments have
% been run. Do not include a kernel-size claim in a submitted draft before the
% table is populated.
%
% \subsubsection{Effect of Kernel Size}
% \label{sec:kernel_size_ablation}
%
% We vary the convolutional kernel size while keeping the training data,
% optimization schedule, block budget, and sparse-attention implementation
% fixed. This experiment tests whether the $7\times7$ receptive field provides
% a better balance between local context modeling and refinement overhead than
% smaller or larger kernels.
%
% \begin{table}[t]
%     \centering
%     \caption{Kernel-size ablation under a matched evaluation protocol.}
%     \label{tab:kernel_size_ablation}
%     \small
%     \begin{tabular}{lrrrr}
%         \toprule
%         Kernel & RULER Avg. & LongBench Avg. & Speedup & Conv. parameters \\
%         \midrule
%         $1\times1$ & TODO & TODO & TODO & TODO \\
%         $3\times3$ & TODO & TODO & TODO & TODO \\
%         $5\times5$ & TODO & TODO & TODO & TODO \\
%         $7\times7$ & TODO & TODO & TODO & TODO \\
%         $9\times9$ & TODO & TODO & TODO & TODO \\
%         \bottomrule
%     \end{tabular}
% \end{table}



% Append the following to the existing Experimental Supplement.
\section{Experimental Supplement}
\label{app:experimental_details}

\subsection{Evaluation Protocol}
\label{app:evaluation_protocol}

LongBench averages use all 16 tasks reported in the main-text tables.
RULER averages use 13 subtasks, and the three-length average weights
32K, 64K, and 128K equally. LongBench top-$k$ ratios are 0.70 for Llama
and 0.75 for Qwen3; the RULER main comparison uses 0.70 for Llama and
0.65 for Qwen3. The training ablation also uses 0.65 on Qwen3.
The Qwen3 RULER main result uses the RULER-trained convolutional checkpoint
evaluated in Table~\ref{tab:ruler_training_ablation}.

Top-$k$ here denotes a retained fraction, not an integer count. Top-$p$
uses a cumulative importance threshold and can yield different retained
counts for different scoring methods. The LongBench reference columns
for full attention, MInference, and FlexPrefill are shared between the
Llama top-$k$ and top-$p$ tables.

\subsection{Efficiency Measurement}
\label{app:efficiency_setup}

The efficiency experiment uses top-$p=0.9$ for Conv and Xattn and the Llama Conv checkpoint
continued from step 4000. Full attention is normalized to 1.00$\times$ at each
context length. The timing harness captures Q/K once and evaluates method
adapters on shared Q/K/V tensors. It times the attention path, including
the selection operations called by each adapter, rather than the complete
model forward pass or autoregressive generation. Table~\ref{tab:topp_speed_values} gives the values plotted in
Figure~\ref{fig:topp_speedup}.

% TODO before submission: add the exact GPU model, number of GPUs, batch size,
% numerical precision, warm-up iterations, timed iterations, software versions,
% and the exact timing repetitions used in the reported run.

\begin{table}[t]
    \centering
    \caption{Measured attention-path speedup over full attention. Conv and Xattn use top-$p=0.9$.}
    \label{tab:topp_speed_values}
    \small
    \resizebox{\linewidth}{!}{%
    \begin{tabular}{lrrrrrr}
        \toprule
        Method & 4K & 8K & 16K & 32K & 64K & 128K \\
        \midrule
        FlexPrefill & 0.07 & 0.12 & 0.35 & 0.72 & 1.30 & 2.65 \\
        Xattn & 0.16 & 0.39 & 0.88 & 1.42 & 2.06 & 2.87 \\
        Conv & 0.13 & 0.35 & 1.04 & 1.76 & 2.61 & 3.67 \\
        MInference & 0.01 & 0.02 & 0.06 & 0.17 & 0.48 & 1.80 \\
        Full attention & 1.00 & 1.00 & 1.00 & 1.00 & 1.00 & 1.00 \\
        \bottomrule
    \end{tabular}%
    }
\end{table}

